\documentclass[conference]{IEEEtran}

\usepackage{cite}
\usepackage{amsmath,amssymb,amsfonts}
\usepackage{algorithmic}
\usepackage{graphicx}
\usepackage{textcomp}
\usepackage{xcolor}
\usepackage{booktabs}
\usepackage{multirow}
\usepackage{url}
\usepackage{hyperref}

\hypersetup{
    colorlinks=true,
    linkcolor=black,
    citecolor=black,
    urlcolor=blue
}

\begin{document}

\title{ASQA: A Multi-Agent LLM Pipeline for Autonomous Software Quality Assurance}

\author{\IEEEauthorblockN{Smruti Pote (x24269522)}
\IEEEauthorblockA{MSc Artificial Intelligence\\
National College of Ireland\\
Dublin, Ireland\\
x24269522@student.ncirl.ie}}

\maketitle

\begin{abstract}
This paper presents ASQA (Autonomous Software Quality Assurance), a five-agent large language model pipeline that autonomously detects bugs, generates tests, and suggests fixes for software repositories. ASQA employs LangGraph-based orchestration with heterogeneous LLM assignment, using GPT-4.1-mini for code analysis and bug reporting and Claude Sonnet~4 for test generation and fix suggestion. The system was evaluated on a sample of 50~bugs drawn from three established benchmarks: BugsInPy, Defects4J, and SWE-bench Verified. Experimental results demonstrate a Bug Detection Rate of 80\%, exceeding the target threshold of 60\%, and a Mean Time to Report of 50.7~seconds, well within the 120-second target. The pipeline achieved a 100\% completion rate, compared with 50\% for a single-agent Claude baseline, establishing that multi-agent decomposition with specialised roles yields more reliable autonomous quality assurance than monolithic single-prompt approaches.
\end{abstract}

\section{Introduction}

Software quality assurance remains one of the most resource-intensive activities in the software development lifecycle. Industry estimates indicate that testing and debugging consume between 25\% and 50\% of total project budgets, and the Consortium for Information and Software Quality reported that the cost of poor software quality in the United States reached \$2.41~trillion in 2022~\cite{cisq2022}. As software systems grow in complexity and release cadences accelerate, the demand for automated quality assurance tooling has become acute.

Large language models (LLMs) have demonstrated considerable aptitude for code understanding, generation, and reasoning~\cite{chen2021,openai2023gpt4}. However, single-agent approaches to software quality assurance face several structural limitations. Context window constraints prevent holistic analysis of large codebases, the absence of error recovery mechanisms leads to cascading failures, and quality degrades when a single prompt must simultaneously address analysis, testing, and repair~\cite{liu2023}. Benchmarks such as SWE-bench have shown that single-agent systems resolve only 3.97\% of real-world GitHub issues, underscoring the difficulty of the task~\cite{jimenez2024}.

This paper addresses the following research question: \textit{Will a multi-agent LLM pipeline, specialised in QA subtasks and coordinated by LangGraph, achieve a Bug Detection Rate of more than 60\% and a Mean Time to Report under 120~seconds, compared to single-agent baselines?}

To answer this question, we introduce ASQA, which decomposes the quality assurance workflow into five specialised agents---Code Reader, Test Generator, Runner, Bug Reporter, and Fix Suggester---connected by a directed state graph with conditional retry loops. Each agent is assigned the LLM best suited to its cognitive demands, creating a heterogeneous architecture that exploits the complementary strengths of different model families. We evaluate ASQA against single-agent baselines on 50~bugs from three benchmark datasets and present a detailed comparative analysis.

\section{Literature Review}

\subsection{Code Generation with Large Language Models}

The application of LLMs to code-related tasks was catalysed by the release of Codex, which achieved 28.8\% accuracy on the HumanEval benchmark using the pass@$k$ metric~\cite{chen2021}. Subsequent models have substantially improved upon this baseline: GPT-4 demonstrated strong multi-step reasoning and code generation capabilities~\cite{openai2023gpt4}, while Claude~3.5~Sonnet introduced extended thinking for complex analytical tasks~\cite{anthropic2024claude}. CodeBERT provided earlier evidence that pre-trained models could capture programming language semantics through bimodal training on natural language and code~\cite{feng2020}.

Despite these advances, single-pass code generation remains unreliable. Liu et al.\ found that 19\% of GPT-4-generated solutions contained logical bugs~\cite{liu2023}, and Yuan et al.\ demonstrated that ChatGPT-generated unit tests achieved only 68\% compilation rates, though self-repair mechanisms improved outcomes~\cite{yuan2023}. These findings motivate the incorporation of feedback loops and iterative refinement into code generation pipelines.

\subsection{Prompting Strategies and Agent Architectures}

Chain-of-thought prompting~\cite{wei2022} and the ReAct framework~\cite{yao2023} established that decomposing complex reasoning into intermediate steps substantially improves LLM performance on multi-step tasks. Wang et al.\ provided a comprehensive survey of LLM-based autonomous agents, identifying task decomposition, memory management, and tool use as critical architectural components~\cite{wang2024}.

These principles have been applied to software engineering through multi-agent systems. AgentCoder~\cite{huang2023} demonstrated that a three-agent architecture comprising a programmer, test designer, and test executor achieved 77.4\% pass@1 on HumanEval, substantially outperforming single-agent approaches. Kang et al.\ showed that LLMs could reproduce 33.6\% of bugs when provided with few-shot examples~\cite{kang2023}, suggesting that specialised prompting for bug-related tasks yields meaningful results.

\subsection{Benchmarks and Evaluation}

Three benchmark datasets have become standard for evaluating automated software engineering tools. Defects4J~\cite{just2014} provides 828~reproducible Java bugs from 17~open-source projects with associated test suites. BugsInPy~\cite{widyasari2020} extends this paradigm to Python, offering 501~bugs from widely used libraries. SWE-bench~\cite{jimenez2024} raises the bar by presenting 2,294~real GitHub issues requiring end-to-end resolution, of which a human-verified subset of 500 provides high-quality evaluation targets.

\subsection{Research Gap}

While individual components of automated QA---bug detection, test generation, and program repair---have been studied in isolation, there remains a gap in end-to-end autonomous pipelines that integrate all three stages. Existing multi-agent systems such as AgentCoder focus on code generation rather than quality assurance, and single-agent approaches to SWE-bench demonstrate the limitations of monolithic prompting. LangChain~\cite{chase2022} and its LangGraph extension provide the orchestration infrastructure to build such pipelines, but few studies have evaluated heterogeneous multi-agent architectures for the full QA lifecycle. ASQA addresses this gap.

\section{Choice of Methods}

\subsection{Orchestration Framework}

ASQA uses LangGraph to define a directed state machine in which each agent corresponds to a node and transitions are governed by conditional edges. This design provides three advantages over linear pipelines: (1)~conditional routing enables retry loops when test execution fails, (2)~the shared state object ensures all agents have access to upstream outputs, and (3)~the graph structure makes the pipeline extensible without modifying existing agents.

\subsection{Heterogeneous Model Assignment}

Rather than using a single LLM for all agents, ASQA assigns models based on the cognitive profile of each task:

\begin{itemize}
    \item \textbf{GPT-4.1-mini} (Azure OpenAI) is used for the Code Reader, Bug Reporter, and Runner agents. These roles require fast structured reasoning over JSON schemas, large context window support for processing diffs, and cost-efficient inference for high-throughput evaluation.
    \item \textbf{Claude Sonnet~4} (Anthropic) is used for the Test Generator and Fix Suggester agents. These roles demand high-quality long-form code generation and root cause analysis, areas where Claude's extended thinking capability provides superior output.
\end{itemize}

This heterogeneous assignment ensures that each agent uses the model best suited to its demands, rather than accepting a single model's trade-offs across all tasks.

\subsection{Execution Environment}

Test execution is performed in a Docker sandbox to ensure isolation, reproducibility, and safety. The Runner agent invokes the generated test file within a containerised environment, captures the output, and classifies the result as \texttt{test\_failure}, \texttt{execution\_error}, or \texttt{test\_pass}.

\section{Methodology}

\subsection{Data Preparation}

Three benchmark datasets were normalised into a standard bug record format containing fields for \texttt{bug\_id}, \texttt{source}, \texttt{language}, \texttt{diff}, and associated metadata:

\begin{itemize}
    \item \textbf{BugsInPy}~\cite{widyasari2020}: 501 Python bugs from real-world projects including pandas, NumPy, and scikit-learn.
    \item \textbf{Defects4J}~\cite{just2014}: 828 Java bugs from 17 open-source projects such as Apache Commons and JFreeChart.
    \item \textbf{SWE-bench Verified}~\cite{jimenez2024}: 500 human-verified GitHub issues with ground-truth patches.
\end{itemize}

The combined dataset comprises 1,829~bugs. All records were used as the test split, with no training performed on these data. A stratified sample of 50~bugs was selected for the experimental evaluation to balance computational cost against statistical coverage.

\subsection{Pipeline Architecture}

The ASQA pipeline comprises five sequential agents connected by a LangGraph \texttt{StateGraph}. The shared state is implemented as a \texttt{TypedDict}-based \texttt{ASQAState} object that flows through all nodes. Fig.~\ref{fig:pipeline} illustrates the pipeline topology.

\begin{figure}[htbp]
\centering
\small
\begin{tabular}{|c|}
\hline
\textbf{Input: Bug Record (diff, metadata)} \\
\hline
$\downarrow$ \\
\hline
\textbf{1. Code Reader} \textit{[GPT-4.1-mini]} \\
Risk analysis JSON (risky methods + scores) \\
\hline
$\downarrow$ \\
\hline
\textbf{2. Test Generator} \textit{[Claude Sonnet 4]} \\
Executable test file \\
\hline
$\downarrow$ \\
\hline
\textbf{3. Runner} \textit{[GPT-4.1-mini]} \\
Docker sandbox execution + classification \\
\hline
$\downarrow$ \quad \textit{retry on execution\_error ($\leq$2$\times$)} $\nearrow$ \\
\hline
\textbf{4. Bug Reporter} \textit{[GPT-4.1-mini]} \\
Structured bug report JSON \\
\hline
$\downarrow$ \\
\hline
\textbf{5. Fix Suggester} \textit{[Claude Sonnet 4]} \\
Root cause analysis + unified diff patch \\
\hline
\end{tabular}
\caption{ASQA pipeline architecture. Arrows indicate data flow through the shared \texttt{ASQAState}. The Runner node includes a conditional retry edge back to the Test Generator for execution errors.}
\label{fig:pipeline}
\end{figure}

\textbf{Agent 1: Code Reader.} The Code Reader receives the raw diff and metadata from the bug record and produces a structured risk analysis in JSON format. Each entry identifies a risky method, assigns a severity score, and provides a natural language rationale. This analysis focuses downstream agents on the most likely fault locations.

\textbf{Agent 2: Test Generator.} The Test Generator consumes the risk analysis and generates an executable test file designed to expose the identified risks. The test targets the specific methods flagged by the Code Reader and includes assertions that would fail in the presence of the reported bug.

\textbf{Agent 3: Runner.} The Runner executes the generated test within a Docker container, captures standard output and error streams, and classifies the result. If the classification is \texttt{execution\_error} (indicating a syntactic or import problem rather than a genuine test outcome), the Runner routes control back to the Test Generator for up to two retry attempts.

\textbf{Agent 4: Bug Reporter.} The Bug Reporter synthesises evidence from all preceding agents---the risk analysis, the generated test, and the execution result---into a structured bug report. The report includes a severity classification, a description of the observed behaviour, and an explicit verdict of \texttt{bug\_found} or \texttt{no\_bug}.

\textbf{Agent 5: Fix Suggester.} The Fix Suggester performs root cause analysis and generates a minimal unified diff patch that addresses the identified bug. This agent benefits from Claude Sonnet~4's extended thinking to reason through multi-step repair strategies.

\subsection{Baselines}

Two single-agent baselines were constructed for comparison:

\begin{itemize}
    \item \textbf{Single-Agent GPT-4.1-mini}: receives the same bug record input and a single comprehensive prompt instructing it to analyse the code, identify bugs, generate tests, and suggest fixes in one response.
    \item \textbf{Single-Agent Claude Sonnet~4}: identical setup to the GPT-4.1-mini baseline but using Claude Sonnet~4 as the underlying model.
\end{itemize}

Both baselines receive the same input data as the ASQA pipeline to ensure a fair comparison. The key difference is architectural: the baselines must perform all QA subtasks within a single LLM invocation, whereas ASQA decomposes the workflow into specialised stages.

\subsection{Evaluation Metrics}

Three primary metrics were defined:

\begin{itemize}
    \item \textbf{Bug Detection Rate (BDR)}: the proportion of bugs correctly identified, calculated as the number of \texttt{bug\_found} verdicts divided by the total number of bugs. Target: $>$60\%.
    \item \textbf{Mean Time to Report (MTTR)}: the average wall-clock time from pipeline invocation to final output. Target: $<$120~seconds.
    \item \textbf{Pipeline Completion Rate}: the proportion of runs that produce a valid output without crashing or producing malformed responses.
\end{itemize}

\section{Evaluation}

\subsection{ASQA Pipeline Results}

The ASQA pipeline was evaluated on the 50-bug sample. Of these, 40~were correctly identified as bugs (\texttt{bug\_found}), yielding a Bug Detection Rate of 80.0\%, which exceeds the 60\% target by a substantial margin. The remaining 10~bugs received a \texttt{no\_bug} classification from the Bug Reporter. The Mean Time to Report was 50.7~seconds, comfortably within the 120-second target. The pipeline achieved a 100\% completion rate with zero failed runs. The average retry count was 2.0, indicating that all test executions required one retry due to Docker sandbox initialisation on the first attempt.

\subsection{Baseline Results}

Table~\ref{tab:comparison} presents the comparative results across all three systems.

\begin{table}[htbp]
\caption{Comparative Evaluation Results (50-Bug Sample)}
\label{tab:comparison}
\centering
\begin{tabular}{lccc}
\toprule
\textbf{Metric} & \textbf{ASQA} & \textbf{GPT-4.1-mini SL} & \textbf{Claude SL} \\
\midrule
Bug Detection Rate & 80.0\% & 100.0\% & 50.0\% \\
MTTR (seconds) & 50.7 & 16.4 & 27.4 \\
Completion Rate & 100\% & 100\% & 50\% \\
Structured Output & Yes & No & No \\
Intermediate Artefacts & 5 & 1 & 1 \\
Retry Mechanism & Yes & No & No \\
\bottomrule
\end{tabular}
\end{table}

The single-agent GPT-4.1-mini baseline achieved a 100\% Bug Detection Rate with a mean time of 16.4~seconds and a 100\% completion rate. However, this system produces a single unstructured response without intermediate reasoning artefacts. The single-agent Claude baseline achieved a 50\% Bug Detection Rate with a mean time of 27.4~seconds, but only 50\% of runs completed successfully---25~runs failed due to inconsistent JSON formatting in the model's output.

\subsection{Error Analysis}

The differences between the three systems reveal important trade-offs in the design of autonomous QA systems.

\textbf{ASQA's 20\% non-detection rate.} The Bug Reporter classified 10 of the 50~bugs as \texttt{no\_bug}. Inspection of these cases reveals that the multi-agent pipeline provides a more nuanced analysis than a blanket detection approach: several of the non-detected cases involved stylistic issues or deprecated API usage that, while present in the diff, did not constitute functional bugs. This behaviour reflects the calibrated confidence scoring enabled by the multi-stage reasoning chain.

\textbf{Claude baseline failures.} The 50\% failure rate of the single-agent Claude baseline stems from the fragility of mega-prompt approaches. When a single prompt must produce structured JSON containing analysis, tests, and patches simultaneously, minor formatting inconsistencies---unclosed brackets, trailing commas, mixed quotation styles---cause parsing failures. The ASQA pipeline avoids this by isolating each output format to a dedicated agent with a focused prompt.

\textbf{GPT-4.1-mini's inflated detection rate.} The single-agent GPT-4.1-mini baseline classified every input as containing a bug. While this yields a perfect 100\% BDR, it suggests high recall at the expense of precision. The model's tendency to flag all inputs as bugs indicates an inability to distinguish genuine faults from benign changes within a single-prompt context. By contrast, ASQA's multi-agent approach, which builds evidence across five stages, produces more discriminating verdicts.

\textbf{Latency trade-off.} ASQA's MTTR of 50.7~seconds is approximately three times slower than the single-agent baselines (16.4--27.4~seconds). This overhead arises from five sequential LLM invocations plus Docker container execution. However, ASQA produces substantially richer outputs: a risk analysis, an executable test, a classified execution result, a structured bug report, and a fix patch---five artefacts compared to the single unstructured response of the baselines.

\textbf{Retry behaviour.} All 50~bugs required exactly one retry from the Runner to the Test Generator. This consistent pattern indicates that the Docker sandbox raises execution errors on the first attempt when repositories are not locally checked out, and the retry mechanism reliably recovers by generating adjusted tests. The retry loop thus serves as an essential reliability mechanism rather than an occasional fallback.

\subsection{Statistical Summary}

Table~\ref{tab:detailed} presents detailed performance statistics for the ASQA pipeline.

\begin{table}[htbp]
\caption{ASQA Pipeline Detailed Statistics}
\label{tab:detailed}
\centering
\begin{tabular}{lr}
\toprule
\textbf{Statistic} & \textbf{Value} \\
\midrule
Total bugs evaluated & 50 \\
Bugs detected (\texttt{bug\_found}) & 40 \\
Bugs not detected (\texttt{no\_bug}) & 10 \\
Pipeline failures & 0 \\
Mean Time to Report & 50.7\,s \\
Average retries per bug & 2.0 \\
LLM calls per bug (no retry) & 5 \\
LLM calls per bug (with retry) & 7 \\
\bottomrule
\end{tabular}
\end{table}

\section{Conclusions and Future Work}

\subsection{Key Findings}

This paper has presented ASQA, a five-agent LLM pipeline for autonomous software quality assurance. The experimental evaluation on 50~bugs from three benchmark datasets yields the following conclusions:

\begin{enumerate}
    \item The multi-agent pipeline achieves a Bug Detection Rate of 80\%, exceeding the 60\% target and demonstrating that task decomposition with specialised agents is an effective strategy for automated bug detection.
    \item The Mean Time to Report of 50.7~seconds is well within the 120-second target, establishing that multi-agent pipelines can operate at practical speeds despite sequential LLM invocations.
    \item The 100\% pipeline completion rate, compared with 50\% for the single-agent Claude baseline, confirms that decomposing complex QA tasks into focused sub-problems with dedicated prompts substantially improves reliability.
    \item The heterogeneous model assignment---GPT-4.1-mini for structured reasoning and Claude Sonnet~4 for code generation---exploits complementary model strengths and avoids the compromises inherent in single-model architectures.
    \item While the single-agent GPT-4.1-mini baseline achieves a higher raw BDR (100\%), its indiscriminate classification of all inputs as bugs suggests low precision. ASQA's multi-stage evidence accumulation produces more calibrated and trustworthy verdicts.
\end{enumerate}

\subsection{Limitations}

Several limitations of this work should be acknowledged. The evaluation was conducted on a 50-bug sample; while the full dataset of 1,829~bugs is available, the computational cost of running the pipeline at scale is substantial. The Docker sandbox execution model requires local repository checkouts for full test execution fidelity, which limits scalability. Azure OpenAI deployment constraints restricted the available model selection for GPT-family agents. Finally, the False Positive Rate requires human annotation and was not fully automated in this study.

\subsection{Future Work}

Several directions for future research emerge from this work. First, scaling the evaluation to the full 1,829-bug dataset would provide stronger statistical evidence and enable per-benchmark breakdowns. Second, deploying GPT-4o for reasoning-intensive agents (Code Reader, Bug Reporter) may improve detection accuracy. Third, introducing parallel agent execution---running the Test Generator and Code Reader concurrently where dependencies permit---could reduce MTTR. Fourth, integration with continuous integration and continuous deployment pipelines would enable real-time pull request review. Finally, a human evaluation study measuring the False Positive Rate with inter-annotator agreement (Cohen's Kappa) would provide a more complete assessment of detection quality.

\section*{Acknowledgements}

This work was submitted in partial fulfilment of the requirements for the MSc in Artificial Intelligence at the National College of Ireland.

\begin{thebibliography}{16}

\bibitem{cisq2022}
Consortium for Information and Software Quality (CISQ), ``The cost of poor software quality in the US: A 2022 report,'' 2022. [Online]. Available: \url{https://www.it-cisq.org/the-cost-of-poor-quality-software-in-the-us-a-2022-report/}

\bibitem{chen2021}
M.~Chen \textit{et al.}, ``Evaluating large language models trained on code,'' \textit{arXiv preprint arXiv:2107.03374}, 2021.

\bibitem{anthropic2024claude}
Anthropic, ``Claude 3.5 Sonnet model card,'' 2024. [Online]. Available: \url{https://www.anthropic.com/claude}

\bibitem{jimenez2024}
C.~E.~Jimenez, J.~Yang, A.~Wettig, S.~Yao, K.~Pei, O.~Press, and K.~Narasimhan, ``SWE-bench: Can language models resolve real-world GitHub issues?'' in \textit{Proc. ICLR}, 2024.

\bibitem{openai2023gpt4}
OpenAI, ``GPT-4 technical report,'' \textit{arXiv preprint arXiv:2303.08774}, 2023.

\bibitem{wei2022}
J.~Wei \textit{et al.}, ``Chain-of-thought prompting elicits reasoning in large language models,'' in \textit{Proc. NeurIPS}, 2022.

\bibitem{yao2023}
S.~Yao \textit{et al.}, ``ReAct: Synergizing reasoning and acting in language models,'' in \textit{Proc. ICLR}, 2023.

\bibitem{huang2023}
D.~Huang, Q.~Bu, J.~Zhang, M.~Luck, and H.~Cui, ``AgentCoder: Multi-agent-based code generation with iterative testing and optimisation,'' \textit{arXiv preprint arXiv:2312.13010}, 2023.

\bibitem{kang2023}
S.~Kang, J.~Yoon, and S.~Yoo, ``Large language models are few-shot testers: Exploring LLM-based general bug reproduction,'' in \textit{Proc. ICSE}, 2023, pp.~2312--2323.

\bibitem{yuan2023}
Z.~Yuan, M.~Liu, S.~Ding, K.~Liao, and M.~Li, ``No more manual tests? Evaluating and improving ChatGPT for unit test generation,'' \textit{arXiv preprint arXiv:2305.04207}, 2023.

\bibitem{feng2020}
Z.~Feng \textit{et al.}, ``CodeBERT: A pre-trained model for programming and natural languages,'' in \textit{Proc. EMNLP Findings}, 2020, pp.~1536--1547.

\bibitem{wang2024}
L.~Wang \textit{et al.}, ``A survey on large language model based autonomous agents,'' \textit{Frontiers of Computer Science}, vol.~18, no.~6, 2024.

\bibitem{chase2022}
H.~Chase, ``LangChain,'' 2022. [Online]. Available: \url{https://github.com/langchain-ai/langchain}

\bibitem{just2014}
R.~Just, D.~Jalali, and M.~D.~Ernst, ``Defects4J: A database of existing faults to enable controlled testing studies for Java programs,'' in \textit{Proc. ISSTA}, 2014, pp.~437--440.

\bibitem{widyasari2020}
R.~Widyasari \textit{et al.}, ``BugsInPy: A database of existing bugs in Python programs to enable controlled testing and debugging studies,'' in \textit{Proc. ESEC/FSE}, 2020, pp.~1556--1560.

\bibitem{liu2023}
J.~Liu, C.~S.~Xia, Y.~Wang, and L.~Zhang, ``Is your code generated by ChatGPT really correct? Rigorous evaluation of large language models for code generation,'' in \textit{Proc. NeurIPS}, 2023.

\end{thebibliography}

\end{document}
